\section{Auxiliary subvarieties over fields of bounded degree}
\label{sec:minimal}

We now treat sequence data with a low factor.  Assume that such data violate
Theorem~\ref{thm:sequence}, and choose a counterexample with
$\dim\cA$ minimal.  The all-high case has already been excluded by
Theorem~\ref{thm:all-high}.  Consequently the projection exclusion of
Lemma~\ref{lemma:projection} and the tangent positivity of
Proposition~\ref{prop:tangent} apply to these data.

Our next objective is to choose smaller varieties through the marked points
without losing control of their equations or their fields of definition.
We fix symmetric very ample line bundles $L_\ell$ on $A_\ell$ and the
resulting projective embeddings.  Write $h_\ell$ for the absolute logarithmic
projective height.  It differs from the corresponding canonical height by
a bounded amount.  All fixed varieties, embeddings, polarizations and bases
of homomorphism groups are defined over one number field $k$, which we
enlarge to contain $\sqrt{-1}$.  Constants in this section may depend on
these fixed data.  They do not depend on a field of definition of a marked
point.

\subsection{Small equations without a degree bound on the point}

\begin{lemma}\label{lemma:height-kernel}
Fix integers $N,D\geq 1$.  Let $F$ be a number field and let $\mathcal S$ be
a set of points in $\mathbb P^N(\Qbar)$ of height at most $T$, where $T\geq1$.
Suppose the vector space of homogeneous degree-$D$ polynomials vanishing
on $\mathcal S$ is defined over $F$.  This space has an $F$-basis whose
coefficient heights are $O_{N,D}(T)$.
\end{lemma}

\begin{proof}
Let $q=\binom{N+D}{D}$ and form the rows of evaluations of the degree-$D$
monomials at points of $\mathcal S$.  Each row is a projective vector of
height at most $DT$.  Select a maximal independent set of these rows;
there are at most $q$ of them.  Their kernel is the required polynomial
space.  Choose pivot columns and give the free variables their standard
basis values.  The resulting kernel vectors have entries which are ratios
of minors of the selected matrix.  The determinant height inequalities
bound their coefficient heights by $O_{N,D}(T)$.

Although the rows can have entries in an extension of arbitrarily large
degree, the resulting vectors lie in $F$.  Indeed the kernel is
$\operatorname{Gal}(\Qbar/F)$-invariant, and its normalized basis with the
chosen free coordinates is unique.  Each automorphism therefore fixes
that basis.  Both the number of rows and all height estimates are
independent of their fields of definition.
\end{proof}

For example, if a degree-$D$ equation over $F$ vanishes at a point $y$, we
apply the lemma to all conjugates of $y$ over $F$.  We obtain a basis of
the equations over $F$ vanishing at $y$, of height $O_D(h(y)+1)$.
If one such equation does not vanish identically on a variety $Z$, at
least one equation in this basis also does not vanish identically on $Z$.

We also need a degree bound for equations that detect smoothness.

\begin{lemma}\label{lemma:bounded-conormal}
Let $Z\subseteq\mathbb P^N$ be an irreducible projective variety of degree
$D$ over an algebraically closed field of characteristic zero.  At every
smooth point of $Z$, the differentials of homogeneous equations of $Z$
of degree at most $D$, dehomogenized in an affine chart, span its
conormal space.
\end{lemma}

\begin{proof}
Linear equations defining the linear span of $Z$ supply its normal
directions in $\mathbb P^N$.  We may therefore work in that span and
assume $Z$ nondegenerate.  Put $d=\dim Z$ and $c=N-d$.  There is nothing
to prove if $c=0$, and the hypersurface equation proves the case $c=1$.

Suppose $c\geq2$ and let $z$ be smooth.  Choose a hyperplane $H$
containing the projective tangent space at $z$.  Since $Z$ is
nondegenerate, $H$ does not contain $Z$.  The closed join of $z$ and
$Z$ has dimension at most $d+1$, and its intersection with $H$ has
dimension at most $d$.  A general $(c-2)$-plane $\Lambda$ in $H$
therefore avoids this join and the tangent space.  Projection from
$\Lambda$ gives a morphism
\[
 \pi:Z\longrightarrow Z'\subseteq\mathbb P^{d+1}.
\]
It is finite because $\pi^*\mathcal O(1)=\mathcal O_Z(1)$ is ample.
It has just one point over $\pi(z)$ and is immersive at $z$.
The scheme-theoretic fiber has length one: its tangent space is zero
and its residue field is the ground field.  For the finite local algebra
over the local ring of $Z'$ at $\pi(z)$, Nakayama's lemma now says that
$1$ generates.  Since $Z'$ is the scheme-theoretic image, the local
map is an isomorphism.  In particular $Z'$ is smooth at $\pi(z)$.

The hypersurface $Z'$ has degree at most $D$.  Its equation pulls back
to an equation of $Z$ whose differential at $z$ defines $H$: the
inverse image of the tangent hyperplane is the span of $\Lambda$ and
the tangent space of $Z$, which is $H$.  Varying $H$ supplies all
conormal directions.
\end{proof}

\subsection{A minimal choice of subproducts}

The choice below adapts the minimal-subproduct descent of
\cite[Section~3]{Remond2005Vojta} and \cite[Section~2]{Dill2020}
to auxiliary varieties whose fields of definition have uniformly bounded
degrees.

Call a sequence of subproducts
\[
 Z_i=\prod_{\ell=0}^s Z_{i\ell},\qquad
 y_{i\ell}\in Z_{i\ell}\subseteq Y_\ell,
\]
\emph{admissible} if the following bounds hold uniformly in $i$:
the geometric degrees of its factors are bounded; all factors are
geometrically integral over a field $k_i\supseteq k$ with
$[k_i:\mathbb Q]$ bounded; and each $Z_{i\ell}$ contains a Zariski dense
set of algebraic points of height at most $C H_{i\ell}$ for one constant
$C$.  We also allow passage to a subsequence in making such a choice.
This definition places no condition on $[k_i(y_{i\ell}):k_i]$.

Admissible products exist: take the fixed parents.  To see the dense
height condition directly, take a finite linear projection of each
parent onto projective space.  The inverse images of the dense set
of projective points with root-of-unity coordinates have bounded
height, by functoriality of heights for that fixed projection.
Since $H_{i\ell}\geq1$, these points give the required bound.

Keep the ambient counterexample fixed.  Among all admissible choices on
all subsequences, minimize the sum of factor dimensions.  This second
minimality concerns the auxiliary products, not the ambient abelian
variety.  After a further extraction, write
\[
 d_\ell=\dim Z_{i\ell},\qquad d=\sum_{\ell=0}^s d_\ell.
\]

\begin{proposition}\label{prop:minimal-products}
For a minimal admissible choice, the following assertions hold after
discarding finitely many indices and passing to a subsequence.
\begin{enumerate}
\item For every pair of positive integers $D,B$, there is an index
$i_0(D,B)$ such that, for every $i\ge i_0(D,B)$ and every factor
$\ell$, any hypersurface of degree at most $D$ defined over an
extension $F_i/k_i$ with $[F_i:k_i]\le B$ which contains
$y_{i\ell}$ also contains $Z_{i\ell}$.
\item For each fixed degree, the equations of $Z_{i\ell}$ in that
degree have a $k_i$-basis of coefficient heights $O(H_{i\ell})$.
The points $y_{i\ell}$ are smooth on $Z_{i\ell}$, and equations of
bounded degree and height $O(H_{i\ell})$ span their conormal spaces.
\item The low factor belongs to a finite set of varieties.  We may
therefore suppose $Z_{i0}=Z_0$ independent of $i$.
\end{enumerate}
\end{proposition}

\begin{proof}
Suppose the first assertion fails for fixed $D,B$.  Choose violating
hypersurfaces along an infinite subsequence, pass to a fixed factor and
a fixed degree, and denote their fields by $F_i$.
Lemma~\ref{lemma:height-kernel}, applied to the
conjugates of $y_{i\ell}$ over $F_i$, replaces each equation by one
of height $O(H_{i\ell})$ which still vanishes at the marked point
and does not vanish identically on $Z_{i\ell}$.  Here we used the
height bound in the original sequence data, rather than any bound
on the degree of the marked point.

Take a geometric irreducible component through $y_{i\ell}$ of the
intersection with this hypersurface.  Its dimension is strictly
smaller than $d_\ell$, and its degree is bounded by ordinary
B\'ezout.  Its conjugates over $F_i$ are also components of this
intersection.  The degree bound therefore bounds the number of
conjugates, and Galois descent defines the chosen component over
an extension of $F_i$ of bounded degree.

We check the dense height condition on that component.  Use the
standard projective adelic metrics, with Fubini--Study metrics at
infinity.  For effective cycles use the additive absolute arithmetic
intersection degree.  The normalized height of an $n$-dimensional
irreducible variety is this degree divided by the product of $n+1$ and
its geometric degree.  The dense height condition, the fundamental height
inequality, and the bounded geometric degrees therefore bound the
arithmetic degree of $Z_{i\ell}$ by $O(H_{i\ell})$.
The arithmetic divisor intersection formula bounds the arithmetic degree
of its intersection with a degree-bounded polynomial by the polynomial
degree times the original arithmetic degree, plus the geometric degree
times the absolute weighted sum of logarithmic sup norms of that
polynomial.  The
latter sum is at most its coefficient height plus a constant
depending only on the degree and the projective space.  Thus the
intersection has arithmetic degree $O(H_{i\ell})$.  Additivity and
nonnegativity give the same bound for the selected component.
Dividing by its positive geometric degree and dimension factor, and
using the other direction of the successive minima inequality, gives a
dense set of points on it of height
$O(H_{i\ell})$; see \cite[Theorem~1.10]{Zhang1995}, or
\cite[Theorem~5.3.3]{YuanZhang2021} with absolute normalization.
The arithmetic divisor estimate used here is the usual
intersection formula for a section, equivalently the
arithmetic divisor formula of
\cite[Proposition~3.2.1(iv), Equation~(3.2.2)]{BostGilletSoule1994};
see also the divisor identity in the proof of
\cite[Corollary~A.4.3]{YuanZhang2021}.
We have produced an admissible product of smaller total dimension,
a contradiction.

For the second assertion apply Lemma~\ref{lemma:height-kernel} to
the dense set of bounded-height points on $Z_{i\ell}$.  Its
evaluation kernel is the space of equations of $Z_{i\ell}$, hence
is defined over $k_i$.  In degrees up to the bounded degree of
$Z_{i\ell}$, these bases span the conormal at every smooth point
by Lemma~\ref{lemma:bounded-conormal}.

Suppose a marked point were singular for infinitely many indices.
Choose an affine chart containing that point.  Among the Jacobian
minors of size equal to the codimension, formed from the bounded
equations just obtained, some minor is nonzero generically on the
variety.  Every such minor vanishes at a singular point: otherwise
the corresponding equations would cut out a smooth germ of the
same dimension containing the germ of the reduced variety, which
would itself then be smooth.  Homogenizing a suitable minor gives
a bounded-degree hypersurface over $k_i$ containing the marked
point and not the variety.  This contradicts the first assertion.
The codimension-zero case has no singular points.

Finally $H_{i0}=1$.  The low-factor equations just constructed have
bounded degrees and bounded coefficient heights over fields of
bounded degree.  Northcott's theorem gives a finite set of possible
coefficient vectors.  The conormal assertion at a smooth point
shows that each $Z_{i0}$ is an irreducible component of a zero locus
of these bounded-degree equations.  There are only finitely many
such zero loci and components.  This proves the last assertion.
\end{proof}

The point of minimality is now explicit: a bounded-degree equation
over a bounded-degree extension which cuts a factor properly is
impossible, regardless of its original coefficient height.  The
remaining argument will construct exactly such an equation.

\subsection{Removing a zero-dimensional low factor}

\begin{lemma}\label{lemma:fixed-low-point}
In a minimal admissible product for the putative low-factor
counterexample, $d_0>0$.
\end{lemma}

\begin{proof}
If $d_0=0$, the fixed low variety is a point $p$, so
$y_{i0}=p$ for every $i$.  Eventual avoidance implies that $p$
lies in no proper torsion coset of $A_0$.

First $V$ projects onto $\Lie(A_0)$.  Otherwise let $P$ be the
orthogonal projector onto the orthogonal complement of that
projection.  This is a nonzero element of $\End(A_0)_{\mathbb R}$.
Writing $\pi_0$ for the low projection, the identity
$P\pi_0=P\pi_0\Phi$ and the small-projection hypothesis give
$P(p)=0$ in the real height space.  But evaluation at $p$ is
injective on $\End(A_0)_{\mathbb Q}$: a nonzero rational
endomorphism taking $p$ to torsion places $p$ in a proper torsion
coset.  Tensoring this injection with $\mathbb R$ proves the
same assertion for real endomorphisms.  This contradicts $P\ne0$.
If there are no high factors, surjectivity already contradicts
the strict codimension hypothesis.

Put $L=\Lie(A_0)$, $H=\bigoplus_{\ell\geq1}\Lie(W_\ell)$,
and write $a=\dim L$, $b=\dim H$, and
$c=\codim_{L\oplus H}V$.  For a positive integer $h$ define
\[
 V_* = \{(v_1,\ldots,v_h)\in H^h:
          \text{there is }v_0\in L\text{ with }(v_0,v_j)\in V
          \text{ for every }j\}.
\]
Kernels and images realize this as the image of a real
homomorphism, hence of a real idempotent.  The space before
forgetting $v_0$ has dimension $a+h(b-c)$, since $V\to L$ is
surjective.  The kernel of forgetting $v_0$ has dimension
$t=\dim(V\cap(L\oplus0))$.  Consequently
\[
 \codim_{H^h}V_* = hc-a+t\geq hc-a.
\]
Since $c>\sum_{\ell=0}^s m_\ell$, a sufficiently large fixed
$h$ makes this strictly greater than
$h\sum_{\ell\geq1}m_\ell$.

Choose $h$ subsequences of indices so that, in the ordered list
of their high factors, the largest scale in one copy divided by
the smallest scale in the next tends to zero.  This is possible
recursively because all the high scales tend to infinity.  Take
the corresponding high marked points, retaining their original
parents and fiber bounds.  All the high nonvanishing hypotheses
remain valid on these subsequences.

Their normalized tuples satisfy the small-projection condition for
$V_*$.  Write $A_+=\prod_{\ell\ge1}A_\ell$ and consider the fixed
real homomorphism whose Lie action is
\[
 T:\Lie(A_0)\oplus\Lie(A_+)^h\longrightarrow\Lie(\cA)^h,
 \qquad T(v_0,v_1,\ldots,v_h)
       =\bigl(\Phi(v_0,v_j)\bigr)_{j=1}^h.
\]
Let $K=\ker T$ and let $\pi_+$ forget the low coordinate.  Since
$\ker\Phi=V\subseteq L\oplus H$, we have $\pi_+(K)=V_*$.
Let $\Phi_*$ be the orthogonal projector with kernel $V_*$ in the
full high ambient space $\Lie(A_+)^h$.  The generalized inverse from
Lemma~\ref{lemma:real-calculus} gives the identity
\[
 \Phi_*\pi_+=\Phi_*\pi_+T^\dagger T,
\]
because $1-T^\dagger T$ projects onto $K$, whose image under
$\pi_+$ is killed by $\Phi_*$.  This identity also holds in the
real height spaces.  At the tuple formed from $p$ and the selected
normalized high points, $T$ tends to zero.  Applying the fixed bounded
operator $\Phi_*\pi_+T^\dagger$ therefore proves the required
small-projection condition.  The high points need only lie in
$\mathcal H(A_+)^h$; the codimension inequality is measured in $H^h$.
We have obtained forbidden all-high sequence data, contrary to
Theorem~\ref{thm:all-high}.
\end{proof}

We henceforth have a fixed positive-dimensional low variety $Z_0$.
The auxiliary varieties in the high factors may still vary.  Their
bounded degrees and fields will suffice for the arithmetic
argument in the next section.
